\section{O Golem de Praga}

No século XVI, a Casa de Habsburgo controlava grande parte da Europa Central, os Países Baixos e a Espanha, bem como as colônias da Espanha nas Américas. A Casa era talvez a primeira verdadeira potência mundial. O Sol brilhava sempre sobre alguma porção dela. Seu governante era também o Imperador do Sacro Império Romano-Germânico, e sua sede de poder era Praga. O Imperador no final do século XVI, Rodolfo II, amava a vida intelectual. Ele investiu nas artes, nas ciências (incluindo astrologia e alquimia) e na matemática, tornando Praga um centro mundial de aprendizado e erudição. É apropriado, então, que nessa atmosfera erudita tenha surgido um robô primitivo, o Golem de Praga.

Um golem (GOH-lem) é um robô de argila do folclore judaico, construído a partir de poeira, fogo e água. Ele é trazido à vida ao inscrever-se \textit{emet}, hebraico para "verdade", em sua testa. Animado pela verdade, mas carente de livre-arbítrio, um golem sempre faz exatamente o que lhe é dito. Isso é uma sorte, porque o golem é incrivelmente poderoso, capaz de suportar e realizar mais do que seus criadores poderiam. No entanto, sua obediência também traz perigo, pois instruções descuidadas ou eventos inesperados podem virar um golem contra seus criadores. Sua abundância de poder é igualada por sua falta de sabedoria.

Em algumas versões da lenda do golem, o Rabino Judah Loew ben Bezalel buscou uma maneira de defender os judeus de Praga. Como em muitas partes da Europa Central do século XVI, os judeus de Praga eram perseguidos. Usando técnicas secretas da \textit{Kabbalah}, o Rabino Judah foi capaz de construir um golem, animá-lo com a "verdade" e ordená-lo a defender o povo judeu de Praga. Nem todos concordaram com a ação de Judah, temendo consequências indesejadas de brincar com o poder da vida. Por fim, Judah foi forçado a destruir o golem, pois sua combinação de poder extraordinário com desajeitamento acabou levando a mortes de inocentes. Apagando uma letra da inscrição \textit{emet} para soletrar, em vez disso, \textit{met}, "morte", o Rabino Judah desativou o robô.

\subsection{Golems estatísticos}

Os cientistas também fazem golems. Nossos golems raramente têm forma física, mas eles também são frequentemente feitos de argila, vivendo no silício como código de computador. Esses golems são modelos científicos. Mas esses golems têm efeitos reais no mundo, por meio das previsões que fazem e das intuições que desafiam ou inspiram. Uma preocupação com a "verdade" anima esses modelos, mas assim como um golem ou um robô moderno, os modelos científicos não são nem verdadeiros nem falsos, nem profetas nem charlatães. Em vez disso, são construtos projetados para algum propósito. Esses construtos são incrivelmente poderosos, conduzindo obedientemente seus cálculos programados.

Às vezes, sua lógica implacável revela implicações anteriormente ocultas para seus criadores. Essas implicações podem ser descobertas inestimáveis. Ou podem produzir comportamentos tolos e perigosos. Em vez de anjos idealizados da razão, os modelos científicos são poderosos robôs de argila sem intenção própria, avançando desajeitadamente de acordo com as instruções míopes que incorporam. Como o golem do Rabino Judah, os golems da ciência são sabiamente vistos com admiração e apreensão. Nós absolutamente temos que usá-los, mas fazê-lo sempre acarreta algum risco.

Existem muitos tipos de modelos estatísticos. Sempre que alguém utiliza até mesmo um procedimento estatístico simples, como um teste $t$ clássico, ela está implantando um pequeno golem que executará obedientemente um cálculo exato, realizando-o da mesma maneira (quase) todas as vezes, sem reclamar. Quase todos os ramos da ciência dependem dos sentidos dos golems estatísticos. Em muitos casos, não é mais possível sequer medir fenômenos de interesse sem fazer uso de um modelo. Para medir a força da seleção natural, a velocidade de um neutrino ou o número de espécies na Amazônia, devemos usar modelos. O golem é uma prótese, fazendo a medição para nós, realizando cálculos impressionantes, encontrando padrões onde nenhum é óbvio.

No entanto, não há sabedoria no golem. Ele não discerne quando o contexto é inadequado para suas respostas. Ele apenas conhece seu próprio procedimento, nada mais. Ele apenas faz o que lhe é dito.

E assim continua sendo um triunfo da ciência estatística que existam agora tantos golems diversos, cada um útil em um contexto particular. Visto dessa forma, a estatística não é matemática nem ciência, mas sim um ramo da engenharia. E, como a engenharia, um conjunto comum de princípios de design e restrições produz uma grande diversidade de aplicações especializadas.

Essa diversidade de aplicações ajuda a explicar por que os cursos introdutórios de estatística são tão frequentemente confusos para os iniciantes. Em vez de um único método para construir, refinar e criticar modelos estatísticos, os alunos recebem um zoológico de golems pré-construídos conhecidos como "testes". Cada teste tem um propósito específico. Árvores de decisão, como a da FIGURA 1.1, são comuns. Respondendo a uma série de perguntas sequenciais, os usuários escolhem o procedimento "correto" para as circunstâncias de sua pesquisa.

Infelizmente, embora estatísticos experientes compreendam a unidade desses procedimentos, estudantes e pesquisadores raramente o fazem. Cursos avançados de estatística enfatizam princípios de engenharia, mas a maioria dos cientistas nunca chega tão longe. Ensinar estatística dessa forma é um pouco como ensinar engenharia de trás para frente, começando com a construção de pontes e terminando com a física básica. Assim, estudantes e muitos cientistas tendem a usar gráficos como a FIGURA 1.1 sem pensar muito em sua estrutura subjacente, sem muita consciência dos modelos que cada procedimento incorpora e sem qualquer estrutura para ajudá-los a fazer os compromissos inevitáveis exigidos pela pesquisa real. Não é culpa deles.

Para alguns, a caixa de ferramentas de golems pré-fabricados é tudo de que precisarão. Desde que permaneçam dentro de contextos bem testados, usando apenas alguns procedimentos diferentes em tarefas apropriadas, muita ciência boa pode ser concluída. Isso é semelhante a como encanadores podem fazer muito trabalho útil sem saber muito sobre dinâmica de fluidos. O problema sério começa quando os estudiosos passam a conduzir pesquisas inovadoras, expandindo as fronteiras de suas especialidades. É como se obtivéssemos nossos engenheiros hidráulicos promovendo encanadores.

Por que os testes não são suficientes para a pesquisa? Os procedimentos clássicos da estatística introdutória tendem a ser inflexíveis e frágeis. Por inflexíveis, quero dizer que eles têm formas muito limitadas de se adaptar a contextos de pesquisa únicos. Por frágeis, quero dizer que eles falham de maneiras imprevisíveis quando aplicados a novos contextos. Isso importa, porque nas fronteiras da maioria das ciências, raramente está claro qual procedimento é apropriado. Nenhum dos golems tradicionais foi avaliado em novos cenários de pesquisa, e por isso pode ser difícil escolher um e depois entender como ele se comporta. Um bom exemplo é o \textit{teste exato de Fisher}, que se aplica (exatamente) a um contexto empírico extremamente estreito, mas é regularmente usado sempre que as contagens de células são pequenas. Eu pessoalmente li centenas de usos do teste exato de Fisher em revistas científicas, mas, além do uso original de Fisher, nunca o vi ser usado adequadamente. Mesmo um procedimento como a regressão linear ordinária, que é bastante flexível de muitas maneiras, sendo capaz de codificar uma grande diversidade de hipóteses interessantes, às vezes é frágil. Por exemplo, se houver um erro substancial de medição nas variáveis de predição, o procedimento pode falhar de maneiras espetaculares. Mas, mais importante, quase sempre é possível fazer melhor do que a regressão linear ordinária, em grande parte devido a um fenômeno conhecido como \textsc{overfitting} (Capítulo 7).

O ponto não é que as ferramentas estatísticas sejam especializadas. Claro que são. O ponto é que as ferramentas clássicas não são diversas o suficiente para lidar com muitas questões comuns de pesquisa. Toda área ativa da ciência lida com dificuldades únicas de medição e interpretação, conversa com teorias idiossincráticas em um dialeto mal compreendido por cientistas de outras tribos. Especialistas em estatística fora da disciplina podem ajudar, mas eles são limitados pela falta de fluência nas preocupações empíricas e teóricas da disciplina.

Além disso, nenhuma ferramenta estatística faz nada por conta própria para abordar o problema básico de inferir causas a partir de evidências. Golems estatísticos não entendem causa e efeito. Eles apenas entendem associação. Sem nossa orientação e ceticismo, golems pré-fabricados podem não fazer nada de útil. Pior, eles podem destruir Praga.

O que os pesquisadores precisam é de alguma teoria unificada de engenharia de golems, um conjunto de princípios para projetar, construir e refinar procedimentos estatísticos de propósito específico. Todo grande ramo da filosofia estatística possui tal teoria unificada. Mas a teoria nunca é ensinada em cursos introdutórios — e frequentemente nem mesmo em avançados. Portanto, há benefícios em repensar a inferência estatística como um conjunto de estratégias, em vez de um conjunto de ferramentas pré-fabricadas.